\section{从 LLVM IR 到 RV64 机器接口}
\label{sec:rv64}

\subsection{合法化、指令选择与寄存器分配}

目标无关代码生成器把 IR 经过调用约定降低、指令合法化（legalization）、指令选择（instruction selection）、寄存器分配和指令调度等步骤，最终打印汇编或发射对象字节\cite{llvm-codegen-2026,llvm-globalisel-2026}。这些步骤是概念分工，不保证每个目标或优化级别都采用相同内部算法。

IR 的 \code{i32 mul} 在 RV64 上要维持 32 位结果语义。作者汇编用 \code{mulw t0,t0,t1} 计算低 32 位并符号扩展到 64 位寄存器；循环变量更新用 \code{addiw}。\code{icmp sle} 与分支则表现为 \code{bgt t1,a0,.Lfactorial\_done}：后端交换比较方向，把“若继续则进循环体”改成“若已结束则跳出”。这是同一控制谓词的目标化布局，不是语义改变。

\begin{lstlisting}[caption={作者 RV64 汇编中的阶乘函数。},label={lst:rvfact}]
factorial:
    li      t0, 1
    li      t1, 2
.Lfactorial_test:
    bgt     t1, a0, .Lfactorial_done
    mulw    t0, t0, t1
    addiw   t1, t1, 1
    j       .Lfactorial_test
.Lfactorial_done:
    mv      a0, t0
    ret
\end{lstlisting}

\code{li}、\code{mv}、\code{j} 和 \code{ret} 是汇编伪指令（pseudo-instruction）：它们改善可读性，由汇编器选成一条或多条真实编码。因此，汇编文本行数不能自动等于对象中的机器指令数。RISC-V ISA 规定真实指令效果，汇编手册和工具约定处理伪指令；本文在编码层以反汇编为准，而非把助记符外观当作硬件事实。

\subsection{RV64 psABI：独立编译仍能相互调用}

LP64D 调用约定规定整数参数和返回值优先使用 \code{a0--a7}；\code{ra} 保存返回地址；\code{sp} 是栈指针并须在过程调用边界满足对齐；\code{t0--t6} 为调用者保存临时寄存器；\code{s0--s11} 为被调用者保存\cite{riscv-psabi-1.0}。因此 \code{factorial} 从 \code{a0} 读入 $n$，在 \code{t0}/\code{t1} 中维护循环状态，再把结果移回 \code{a0}。

叶函数 \code{factorial} 不调用别的函数，可直接 \code{ret}；\code{main} 连续调用五个函数，必须保护自己的返回地址。作者汇编把栈指针减 16，在 \code{8(sp)} 保存 \code{ra}，返回前逆向恢复。这 16 字节而非任意大小来自调用边界的栈对齐约束。它也说明“有局部变量”不必然意味着“有栈槽”：这里寄存器足以承载局部状态，栈帧的主要任务是保存返回地址。

\subsection{结构观测与不能推出的结论}

同一 Clang 18.1.3 环境中，生成汇编的近似指令行由 \code{-O0} 的 77 降到 \code{-O2} 的 45，出现的助记符种类由 13 增至 15。前一变化可能来自栈访问和冗余搬运减少，后一变化说明优化会选用更多样的操作；“种类更多”不等于程序更复杂或更慢。近似计数排除标签和大多数指示符，却仍包含伪指令，因此只适于描述文本结构。

\begin{table}[tb]
\caption{RV64 汇编的结构观测，不是性能测量。}
\centering
\begin{tabular}{@{}lrrl@{}}
\toprule
级别 & 近似指令行 & 助记符种类 & 可支持的判断\\
\midrule
\code{-O0} & 77 & 13 & 保留较多显式搬运/帧操作\\
\code{-O2} & 45 & 15 & 文本更短但指令选择更多样\\
\bottomrule
\end{tabular}
\end{table}

到汇编层，类型信息已大幅弱化：\code{a0} 本身不携带“这是 SysY int”的标签；\code{mulw} 的宽度及调用点约定共同保存所需行为。符号名仍没有最终地址，而 \code{call getint} 甚至尚未决定具体位移。汇编器的任务正是把这些文本承诺编码成节、符号和重定位。
